\section{Results}
\label{sec:results}

\subsection{Evaluation Protocol}

Our primary measurement is macro-scroll Dice on a frozen 689-row benchmark
drawn from the two held-out scrolls: the unweighted mean of the per-scroll
pooled Dice, so a scroll contributes equally regardless of how many rows it
supplies.
The benchmark manifest, its row set, and its threshold grid were fixed before
this study and are hash-checked at evaluation time.

Three protocol choices matter for interpreting what follows.
First, the operating threshold is declared rather than fitted per candidate:
$t=\fzPrimaryT$ is the primary endpoint and $t=\fzShipT$ the single
preregistered secondary point, chosen a priori because it is the calibrated
argmax at 21 of 25 milestones.
A per-candidate calibrated threshold is reported as a diagnostic only, since
selecting it inside the same rows inflates the score.
Second, uncertainty is measured by paired cluster bootstrap: sampling units are
support-anchor chunks within a scroll rather than individual rows, because many
validation rows share an anchor, and a candidate and its reference are resampled
with identical weights so their difference is paired.
Such an interval describes patch sampling and says nothing about seed noise,
which we measure separately.
Third, a promoted candidate must additionally pass a blind six-cube morphology
gate on PHerc1447 and a frontier veto: its reference-relative bridge pixels may
not exceed $1.20$ times the reference curve interpolated at matched skeleton
coverage, with no extrapolation outside the measured coverage overlap.

\subsection{Milestone Behaviour and Seed Noise}

All 25 milestones completed their frozen audits and were gate-eligible at
$t=\fzPrimaryT$.
The response is not a curve that rises and plateaus.
Over the eleven late milestones at or beyond $150{,}000$ exposures, the mean
fixed-threshold macro Dice is $\fzLateMean$ with a standard deviation of
$\fzLateSD$, and adjacent milestones swing by several times that deviation in
both directions.
Batch three with momentum $0.99$ random-walks in a flat basin.

The preregistered late-selection rule chooses \fzSelected\ exposures at
$\fzMacroFull$.
For context rather than selection, the global maximum over all milestones is
$\fzGlobalMax$ and the exact endpoint is $\fzEndpoint$, so a rule that simply
took the last checkpoint would have lost roughly a standard deviation of
milestone noise.

We ran a second full-horizon replicate under an identical recipe with a
different seed in order to know how much of that spread is not attributable to
the method at all.
Its root-mean-square difference from the primary run across the 25 matched
milestones is $\seedRMS$, falling to $\seedRMSLate$ over the late band.
This number is the reason we do not claim duration superiority: the gap between
our selected milestone and an early milestone of the same recipe is smaller
than the difference two seeds of that recipe produce.
It is also the yardstick every ablation in Appendix~\ref{sec:also-tried} is
measured against.

\subsection{Comparison with Released M7}

Table~\ref{tab:headline} reports the frozen benchmark under matched inference.
Against released M7 evaluated identically and at the same threshold, the
student gains $\fzGain$ macro Dice, a 22.9\% relative improvement.
The comparison is not an ensemble: inference runs the raw student alone, with
no teacher forward pass and no probability blend.

\begin{table}[t]
\caption{Frozen 689-row benchmark, macro-scroll Dice. TTA is eight-way mirror
test-time augmentation on unchanged weights. The student's $t$ values are
preregistered; M7's calibrated $t=0.20$ is fitted in-sample and therefore
favours the baseline.}
\label{tab:headline}
\centering
\small
\begin{tabular}{llr}
\toprule
Model & Inference & Macro Dice \\
\midrule
Released M7 & $t=\fzPrimaryT$ & $\mSevenMacro$ \\
Released M7 & TTA, $t=\fzShipT$ & $\mSevenTtaMacro$ \\
Released M7 & TTA, calibrated $t=0.20$ & $\mSevenTtaCalibrated$ \\
\midrule
Ours (\fzSelectedShort) & $t=\fzPrimaryT$ & $\fzMacro$ \\
Ours (\fzSelectedShort) & TTA, $t=\fzPrimaryT$ & $\fzTtaPrimary$ \\
Ours (\fzSelectedShort) & TTA, $t=\fzShipT$ & $\mathbf{\fzTtaMacro}$ \\
\bottomrule
\end{tabular}
\end{table}

\subsection{Head-to-Head at the Published Setting}
\label{sec:head-to-head}

The comparison above runs M7 the way we run the student. The organisers
publish M7 differently: their inference metadata disables test-time
augmentation and thresholds at $0.2$. Table~\ref{tab:h2h} scores that exact
setting and HercUNet v0 \cite{lomeli2026hercunet}, a concurrent iterative
refiner warm-started from M7. Nothing is retrained, every model sees the same
192\textsuperscript{3} patches, and one metric binary computes every composite.

HercUNet publishes no operating threshold, so it receives the best of both
normalisations its code contains, one to four refinement passes, augmentation
on or off, every threshold from $0.10$ to $0.90$ and up to three voxels of
dilation, each chosen on this benchmark; that choice flatters it. Because it
predicts the medial surface of a sheet rather than its face, we also report
surface F1 with two- and four-voxel tolerance. Its own command-line pipeline,
run on the live scroll around four benchmark patches with its blending and
half-window shifts, read CT identical to our patches and scored at or below our
harness on every one.

\begin{table}[t]
\caption{Frozen 689-row benchmark, same rows for every model. Comp.\ is the
Kaggle-style $0.35$ SurfaceDice@2 $+\,0.35$ VOI $+\,0.30$ TopoScore. Surface F1
counts a prediction within the stated distance of a labelled sheet as correct.
HercUNet's row is the best over all of its settings on this benchmark.}
\label{tab:h2h}
\centering
\small
\setlength{\tabcolsep}{4pt}
\begin{tabular}{llrrrr}
\toprule
Model & Inference & Dice & Comp. & F1@2 & F1@4 \\
\midrule
M7, as published & no TTA, $t=0.20$ & $\mSevenPublishedMacro$ & $\mSevenPublishedComposite$ & $\mSevenPublishedTolTwo$ & $\mSevenPublishedTolFour$ \\
HercUNet v0 & best setting & $\hercBestMacro$ & $\hercBestComposite$ & $\hercTolTwo$ & $\hercTolFour$ \\
Ours & TTA, $t=\fzShipT$ & $\mathbf{\fzTtaMacro}$ & $\mathbf{\fzCompositeTta}$ & $\mathbf{\fzTolTwo}$ & $\mathbf{\fzTolFour}$ \\
\bottomrule
\end{tabular}
\end{table}

Paired over the same rows, the student leads M7 as published by
$\fzMinusMSevenPublished$ $[\fzMinusMSevenPublishedLow,\fzMinusMSevenPublishedHigh]$
macro Dice and HercUNet at its best by $\fzMinusHerc$
$[\fzMinusHercLow,\fzMinusHercHigh]$. Dilating HercUNet's medial prediction,
the remedy its own write-up proposes, raises it only to $\hercBestDilatedMacro$.
Its additional refinement passes lower its Dice on this benchmark but, with
augmentation, raise its topology score, so its best composite is the four-pass
output with TTA.

\subsection{Where the Published M7 Makes Blobs}
\label{sec:blob-audit}

The benchmark's labelled regions are relatively clean, so it measures accuracy
more than blobs. To test the blob claim without choosing examples, we drew 32
cubes of 128\textsuperscript{3} voxels uniformly at random from the interior of
each of eight scrolls that carry an official published M7 prediction, 256 in
all, under a rule written into the sampling code before any cube was read. For
each cube we kept the published mask unchanged and ran the student and
HercUNet on the same CT. We measure the share of predicted voxels at
six-connected depth five or more from background, which a single sheet at this
resolution does not reach, and split each scroll's cubes into thirds by how much
of their CT is dark gap between wraps (Table~\ref{tab:blob}).

\begin{table}[t]
\caption{Share of predicted voxels inside a solid mass (depth $\geq 5$), in
percent, over 256 random interior cubes from eight scrolls, split within each
scroll by compression.}
\label{tab:blob}
\centering
\small
\setlength{\tabcolsep}{4pt}
\begin{tabular}{lrrr}
\toprule
Source & open & middle & compressed \\
\midrule
M7, as published & 2.0 & 2.7 & 6.0 \\
M7, released weights rerun, same setting & 2.1 & 2.5 & 6.5 \\
M7 + TTA, $t=0.30$ & 0.4 & 0.2 & 0.6 \\
HercUNet v0, best setting (four passes + TTA) & 5.3 & 5.1 & 7.4 \\
HercUNet v0, four passes, no TTA & 5.0 & 4.8 & 6.6 \\
HercUNet v0, one pass, no TTA, $t=0.50$ & 0.8 & 0.8 & 1.0 \\
Ours, shipped & 2.9 & 1.5 & 2.8 \\
\bottomrule
\end{tabular}
\end{table}

The published mask's solid-mass share triples from open to compressed cubes,
and rerunning the released weights at the same setting reproduces it, so the
effect belongs to the model rather than to the publication pipeline. The
student does not rise with compression. Augmenting M7 removes most deep voxels,
but in compressed cubes what remains is fragments rather than sheets: it fills
less of the cube than the student (17.5\% against 23.0\%) yet breaks into more
pieces ($33.9$ per cube against $24.0$), and on the benchmark that setting still
trails the student by $0.120$ Dice. Whole-cube failures are rare for every model: one of
256 published M7 cubes is a solid slab under the ScrollFiesta rejection rule,
and none of the student's. At its best challenge-metric setting, four passes
with TTA, HercUNet is the thickest output measured.

\subsection{Inference-Time Levers}

Two levers that change no weights were preregistered before scoring, together
with the rule that a claim above $0.70$ requires a positive paired lower bound
after Holm correction and non-negative deltas on both scrolls.
Eight-way mirror TTA at $t=\fzShipT$ reaches $\fzTtaMacro$
$[\fzTtaLow,\fzTtaHigh]$, a paired gain of $\fzTtaDeltaOp$
$[\fzTtaDeltaOpLow,\fzTtaDeltaOpHigh]$ over the same weights without TTA at
their training-record threshold, positive on both scrolls.
At the primary endpoint $t=\fzPrimaryT$ the same weights reach $\fzTtaPrimary$,
a paired gain of $\fzTtaDelta$ $[\fzTtaDeltaLow,\fzTtaDeltaHigh]$ at matched
threshold.
The preregistered replication applied the same recipe to the independent
seed-1204 run: it reaches $\cThreeRTtaMacro$ $[\cThreeRTtaLow,\cThreeRTtaHigh]$,
$\cThreeRDelta$ $[\cThreeRDeltaLow,\cThreeRDeltaHigh]$ over that run's own
no-TTA reference, so the gain is not a property of one seed.

We also evaluated uniform weight averages over late milestones.
They reduce variance without shifting the level: no narrow late average beat the
selected milestone, and only a wide average spanning $100{,}000$ to
$\fzHorizon$ exposures edged past it, by less than the seed band.
We report those variants and do not promote them.

At the shipped operating point the blind six-cube gates and the frontier veto
pass: foreground ratio $0.799$ against the reference, with $4{,}945$
reference-relative bridge pixels against a $7{,}983$ ceiling at the measured
skeleton coverage of $0.782$.
A Kaggle-style composite on the same rows rises from $\fzComposite$ to
$\fzCompositeTta$, with the topology component rising from $\fzTopo$ to
$\fzTopoTta$ and thereby meeting the $\fzTopoAdvisory$ advisory that the
no-TTA configuration missed.

\subsection{Held-Out Human Labels}

On PHerc0500P2, the one held-out set with human rather than model-derived
labels, the composite is flat across models: the shipped student scores
$0.4917$, M7 as published $0.4938$, M7 with TTA at its calibrated threshold
$0.5015$, and HercUNet v0 at most $0.4795$. Every variant of our model we
measured falls between $0.476$ and $0.500$, a spread smaller than the
differences we are trying to resolve.
These are absolute scores only and not a gain claim, because released M7 saw
those labels during its own training and HercUNet starts from M7; we report
them because a flat result on the only human-labelled holdout is a real limit
on what the frozen benchmark establishes.

\subsection{Release}
\label{sec:release}

We release the selected \fzSelectedShort\ checkpoint with its shipped
inference recipe as \releaseName; our records call it F0.
The code and the \checkpointMB\,MB checkpoint are released under the Apache
License 2.0, and the checkpoint is checked against its recorded SHA-256 before
use.
One command runs the shipped recipe, eight-way mirror TTA at $t=\fzShipT$, on
any region of a registered scroll, streaming the CT from the public Vesuvius
Challenge data bucket, and writes an OME-NGFF prediction store together with a
cube grid that downstream geometry tools read directly.
The postprocessors of Appendix~\ref{sec:fitter} ship with it as opt-in tools
pinned to $t=\fzPrimaryT$ (Sec.~\ref{sec:limitations}).
The upstream data, the M7 weights the student was initialised from, and the fine
teacher keep their own terms.

\subsection{Figure Provenance}
\label{sec:figure-provenance}

The teaser (Fig.~\ref{fig:teaser}) shows the shipped model at its operating
point, beside the published M7 mask and HercUNet v0 on seeded-random cubes.
The other rendered panels---Fig.~\ref{fig:registered},
Fig.~\ref{fig:limitations}, and both galleries in
Sec.~\ref{sec:gallery}---show the earlier v31 student at $t=0.45$, not the
model reported in Table~\ref{tab:headline}.
Those panels were produced for the earlier configuration and have not yet been
regenerated.
They remain informative about the qualitative failure modes the method
addresses, and the registered comparison grammar is unchanged, but they should
not be read as depicting the shipped weights.
Regenerating the panel set against the selected checkpoint is outstanding
release work and is listed in Sec.~\ref{sec:limitations}.
The quantitative gate review of the shipped model, including its six-cube
panels, is complete and lives in the release record.

\begin{figure*}[t]
  \centering
  \figureasset{figure_2_registered_examples}{0.96\linewidth}{2.2in}
  \caption{Locked PHerc0139 comparisons at a fixed operating point.
  Columns show the same CT crop, released M7, projected fine-teacher target, raw
  student, and a two-sided teacher comparison. Green denotes shared foreground,
  yellow teacher-only foreground, and magenta student-only foreground.
  Student panels show the v31 configuration; see
  Sec.~\ref{sec:figure-provenance}.}
  \Description{Four registered locked examples comparing CT, released M7,
  projected teacher, raw student, and two-sided difference.}
  \label{fig:registered}
\end{figure*}

\subsection{Interpretation}

The result we are most confident in is a level: under matched inference the
student is far enough above released M7 that no threshold, seed, or milestone
choice within our measured spread reverses the ordering.
The results we hold loosely are the ones inside that spread.
Duration beyond the late band, the choice among late milestones, and the
weight-average variants all sit within seed noise and are reported as
unresolved rather than settled.
Both gate corpora remain model-derived rather than ground truth, so the
morphology review constrains the selected checkpoint without certifying it,
and the flat human-labelled holdout keeps that caveat in view.
